\documentclass[11pt]{article}
\input{../../shared/project_style.tex}
\rhead{Basics 5 Textbook Draft}

\begin{document}
\DocTitle{Basics 5: Partial Differential Equations}{I - Mathematical and computational foundations}{Classification, data requirements, characteristics, conservation form, and well-posedness}

\begin{overviewbox}
\textbf{Textbook draft, version 1.0.} Most plasma simulations solve partial differential equations. Their mathematical type predicts information propagation, boundary data, regularity, and suitable numerical methods. This chapter is designed for a reader with no prior plasma-physics training. It distinguishes exact identities from physical approximations and connects the central equations to later simulation modules.
\end{overviewbox}

\ProjectTOC

\section{Learning objectives}
After completing this note, the reader should be able to:
\begin{itemize}[leftmargin=1.6em]
\item Identify independent variables, dependent fields, order, and linearity.
\item Classify canonical elliptic, parabolic, and hyperbolic equations.
\item Choose appropriate initial and boundary conditions.
\item Derive characteristics for advection.
\item Distinguish conservative and nonconservative forms.
\item Explain weak solutions, eigenmodes, and well-posedness.
\end{itemize}

\section{Prerequisites and reading strategy}
\begin{itemize}[leftmargin=1.6em]
\item Basics 1-3 or equivalent multivariable calculus.
\end{itemize}
On a first reading, emphasize physical meaning and dimensional checks. On a second reading, reproduce the derivations. Attempt the computational laboratory only after the limiting cases are understood.

\section{Definitions and examples}

A PDE relates a field to its partial derivatives:
\begin{equation}
\mathcal F(\mathbf x,t,u,\partial_tu,\nabla u,\nabla^2u,\ldots)=0.
\end{equation}
The order is the highest derivative. The equation is linear if $u$ and its derivatives enter linearly with coefficients independent of $u$. Poisson, diffusion, wave, ideal-MHD, and kinetic equations illustrate different mathematical structures.


\section{Elliptic, parabolic, and hyperbolic types}

For
\begin{equation}
Au_{xx}+2Bu_{xy}+Cu_{yy}+\cdots=0,
\end{equation}
the sign of $B^2-AC$ gives the local classification. Poisson equation is elliptic and globally constrained by boundaries. Diffusion is parabolic and smooths gradients. The wave equation is hyperbolic and carries finite-speed signals. Ideal MHD is a nonlinear hyperbolic evolution system, while equilibrium calculations are closer to elliptic boundary-value problems.


\section{Initial and boundary conditions}

A first-order-in-time equation generally needs one initial field. A second-order spatial operator needs boundary conditions. Common forms are
\begin{align}
u&=g &&\text{Dirichlet},\\
\pdv{u}{n}&=h &&\text{Neumann},\\
au+b\pdv{u}{n}&=c &&\text{Robin}.
\end{align}
Data must be compatible with the operator. Pure-Neumann Poisson problems need an integral compatibility condition and have an arbitrary additive constant.

A problem is well posed if a solution exists, is unique, and depends continuously on data. Numerical sophistication cannot repair a fundamentally ill-posed model.


\section{Characteristics and causality}

For advection,
\begin{equation}
u_t+au_x=0,
\end{equation}
follow a curve $x(t)$:
\begin{equation}
\dv{u}{t}=u_t+\dot xu_x.
\end{equation}
Choosing $\dot x=a$ gives $du/dt=0$, so
\begin{equation}
u(x,t)=u_0(x-at).
\end{equation}
Characteristics identify causal directions and inflow boundaries. For $a>0$, the left boundary is inflow and normally requires data; the right boundary is outflow.


\section{Conservative form and discontinuities}

A scalar conservation law is
\begin{equation}
u_t+\pdv{}{x}F(u)=0.
\end{equation}
For smooth $u$, it is equivalent to $u_t+F'(u)u_x=0$. Across a discontinuity, the conservative form determines shock speed:
\begin{equation}
s=\frac{F(u_R)-F(u_L)}{u_R-u_L}.
\end{equation}
MHD supports shocks, so flux-form discretization is essential. Expanding products before discretization can destroy the correct jump condition.


\section{Separation of variables, spectra, and weak forms}

For the fixed-end wave equation, $u=X(x)T(t)$ gives
\begin{equation}
\frac{T''}{c^2T}=\frac{X''}{X}=-k^2,
\end{equation}
and boundary conditions select $k_n=n\pi/L$. This is the prototype of a normal-mode eigenproblem.

When solutions are not classically differentiable, multiply by a test function and integrate to obtain a weak form. Finite elements are based on weak forms. Standard PINNs often minimize strong residuals at collocation points, which may be inappropriate near shocks or ideal continuum singularities.


\section{Computational laboratory}

Implement one-dimensional Poisson, diffusion, and advection problems on a common grid. Use a direct boundary-value solve, explicit and implicit diffusion stepping, and upwind advection. Demonstrate global boundary influence for Poisson, a stability limit for explicit diffusion, and oscillations from centered advection. Measure convergence against analytic solutions.


\section{Summary}
\begin{itemize}[leftmargin=1.6em]
\item PDE type predicts information propagation and required data.
\item Elliptic, parabolic, and hyperbolic equations behave qualitatively differently.
\item Characteristics encode causality.
\item Conservative form is essential for discontinuous solutions.
\item Weak and strong residual formulations are not interchangeable in all regimes.
\end{itemize}

\SymbolsAcronymsSection
\begin{symtable}
$u(\mathbf x,t)$ & field & Unknown dependent variable. \\
$\mathcal F$ & PDE operator & Relation among field and derivatives. \\
$F(u)$ & flux & Transported quantity in a conservation law. \\
$s$ & shock speed & Discontinuity propagation speed. \\
$\partial_nu$ & normal derivative & $\nabla u\cdot\mathbf n$. \\
\end{symtable}

\section{Exercises}
\begin{enumerate}[leftmargin=1.8em]
\item Classify $u_{xx}+4u_{xy}+u_{yy}=0$.
\item Derive characteristics for $u_t+a(x,t)u_x=0$.
\item Find the compatibility condition for a pure-Neumann Poisson problem.
\item Derive Burgers shock speed.
\item Explain why a wave equation needs two initial conditions while advection needs one.
\end{enumerate}

\section{References}
\begin{thebibliography}{99}
\bibitem{ref1} G. B. Arfken, H. J. Weber, and F. E. Harris, \emph{Mathematical Methods for Physicists}, 7th ed., Academic Press (2013).
\bibitem{ref2} G. Strang, \emph{Linear Algebra and Its Applications}, 4th ed., Brooks/Cole (2006).
\bibitem{ref3} W. A. Strauss, \emph{Partial Differential Equations: An Introduction}, 2nd ed., Wiley (2007).
\bibitem{ref4} J. P. Freidberg, \emph{Ideal MHD}, Cambridge University Press (2014).
\bibitem{ref5} J. P. Goedbloed, R. Keppens, and S. Poedts, \emph{Advanced Magnetohydrodynamics}, Cambridge University Press (2010).
\bibitem{ref6} H. Goedbloed and S. Poedts, \emph{Principles of Magnetohydrodynamics}, Cambridge University Press (2004).
\end{thebibliography}

\end{document}
